% 第 3 章 User Interface and Data Preparation（原书 3-9 … 3-32）
\bichapter{用户界面与数据准备}{User Interface and Data Preparation}
\label{chap:ui}

\section{引言}
\subsection*{Introduction}

本章介绍控制 RedHawk 的 TCL 命令行用户界面与图形用户界面（GUI）、管理输入/输出数据的推荐目录组织方式，以及启动 RedHawk 分析前所需的输入数据文件与准备工作。软件安装说明见附录~\ref{chap:install}。

\section{TCL 命令用户界面}
\subsection*{TCL Command User Interface}

\subsection{TCL 命令摘要}
\subsubsection*{TCL Command Summary}

RedHawk 程序可通过 TCL 命令行界面及本节所述命令集调用与控制。

基本 RedHawk 启动方式为：
\begin{lstlisting}
redhawk
\end{lstlisting}

以下为运行 RedHawk 可用 TCL 命令摘要。各命令的选项与语法详见附录~\ref{chap:cmd}「TCL / Script Commands」。

\begin{itemize}
  \item \cmd{cell swap}——允许将高功耗单元移至压降较低的位置
  \item \cmd{characterize}——运行 Apache Power Library（APL）表征，生成动态 Vdd/Vss 电流波形及 intrinsic/intentional decap 值
  \item \cmd{condition}——设置、显示或取消 DvD 后处理分析的各种条件（如 plot、print）
  \item \cmd{config}——设置 GUI 选项，用于查看、分析与调试 RedHawk 结果
  \item \cmd{decap}——允许对去耦电容进行各类修改，以降低动态压降
  \item \cmd{dump}——以 GIF 格式写出色阶图（colormap）
  \item \cmd{eco}——定义对设计数据库的修改（仅 GUI 操作），等价于 GUI what-if 操作
  \item \cmd{export}——从 RedHawk 引擎导出数据供其他工具使用
  \item \cmd{fao}——允许修改物理设计
  \item \cmd{generate}——生成指定文件供查看和/或后续使用
  \item \cmd{get}——允许查询数据库，按模式匹配搜索并查找单元或实例名、执行脚本，或更深入地访问设计
  \item \cmd{gsr}——获取 GSR 变量值、设置 GSR 值、显示内容，或将 GSR 内容发送到文件
  \item \cmd{history}——显示本会话期间在命令行输入的所有先前命令
  \item \cmd{import}——导入指定文件用于 RedHawk 分析，例如 APL、AVM、DB、DEF、ECO、GSR、guiconf、LEF、LIB、PAD、POWER 或 TECH
  \item \cmd{ircx2tech}——将输入 iRCX 文件转换为适当的 Apache tech 文件
  \item \cmd{marker}——按 x,y 坐标或实例名添加或删除设计标记
  \item \cmd{mesh}——对电源网格进行各类修改（添加、删除、修改宽度与间距），以降低压降
  \item \cmd{message}——显示 Error、Info 与 Warning 消息信息
  \item \cmd{movie}——设置或播放基于实例的动态压降历史「影片」
  \item \cmd{perform}——运行选定的 RedHawk 分析，例如提取、功耗计算、静态分析、动态分析及时序 slew
  \item \cmd{pfs}——执行各类 ESD 分析功能
  \item \cmd{plot}——按指定条件生成图形曲线
  \item \cmd{print}——按指定条件打印基于文本的报告
  \item \cmd{probe}——在提取前选择或取消选择特定节点
  \item \cmd{psiwinder}——对设计执行各类时钟与时序分析
  \item \cmd{query}——显示所选设计对象的信息与分析参数
  \item \cmd{report}——创建 RedHawk 分析报告
  \item \cmd{ring}——添加或删除电源环（power ring）
  \item \cmd{route fix}——布线新电源/地线并添加相关过孔，以降低过大压降
  \item \cmd{save}——保存设计数据库供后续使用
  \item \cmd{select}——在 GUI 中选择并高亮对象
  \item \cmd{setup}——设置执行 RedHawk 分析所需的数据与条件
  \item \cmd{show}——在 GUI 中显示指定 RedHawk 结果的色阶图
  \item \cmd{window}——更改 GUI 窗口几何，使 \cmd{dump gif} 命令可生成更高分辨率色阶图
  \item \cmd{zoom}——放大或缩小设计视图
\end{itemize}

RedHawk 图形用户界面及其功能见下一节。

\section{图形用户界面}
\subsection*{Graphical User Interface}

\subsection{GUI 组成元素}
\subsubsection*{Elements of the GUI}

图形用户界面便于查看设计元素与 RedHawk 分析结果。从命令行调用 RedHawk 时，GUI 会自动显示。

RedHawk GUI 包含多个功能面板（图~\ref{fig:rh-gui}），主要功能区域为：
\begin{itemize}
  \item 菜单栏（Menu bar）
  \item 主显示区（Primary display area）
  \item 多页视图标签（Multiple page view tabs）
  \item 文本显示区（Text display area）
  \item 控制按钮（Control buttons）
  \item 设计视图区（Design view area）
  \item 光标位置读数（Cursor location readout）
  \item TCL 命令行（TCL command line）
\end{itemize}

以下各节概述这些 GUI 元素的功能。GUI 命令与函数的详细说明见附录~\ref{chap:cmd}「RedHawk Graphic User Interface Description」。

\figplaceholder{Figure 3-1}{RedHawk 图形用户界面}
\label{fig:rh-gui}

\subsection{菜单栏}
\subsubsection*{Menu Bar}

菜单栏可用的 RedHawk 功能如下：
\begin{itemize}
  \item \textbf{File}——导入与导出设计数据、数据库、ECO 信息及配置信息
  \item \textbf{Edit}——编辑 GSR 与 Tech 文件中的信息，以及修改电源网格元素（如 pad、strap、via）
  \item \textbf{View}——显示并更改色阶图，查看设计元素（如网表、层、实例、电容及一般对象属性）与分析结果（如 IR 与动态压降、功耗与电流密度、EM 问题、slack 与 delay）
  \item \textbf{Tools}——调用 RedHawk 工具，例如电源网格与去耦电容修复优化、低功耗电路设计分析，以及 PJX 关键路径与时钟树分析
  \item \textbf{Static}——执行与静态压降分析相关的功能，例如功耗计算、网络 R 提取、pad 与封装约束，以及 IR drop 与 EM 分析
  \item \textbf{Dynamic}——执行与动态压降分析相关的功能，例如功耗计算、网络 RLC 提取、pad 与封装约束，以及 vectorless 或基于 VCD 的动态压降分析
  \item \textbf{Timing}——执行与电路时序分析相关的功能，例如实例 slew、slack、K-factor、delay、时钟网、修改 SDF 数据，以及压降对时序的影响
  \item \textbf{Results}——设置并生成 RedHawk 分析结果信息，例如直方图、日志消息，以及按最差 IR drop、动态压降、电迁移、电源/地设计薄弱点排序的实例或节点列表；还可创建并显示按时间步显示的压降「影片」
  \item \textbf{Explorer}——控制分析结果的生成与显示
  \item \textbf{Windows}——设置并创建多页同时显示，以及设置窗口显示首选项
  \item \textbf{Help}——显示 RedHawk 软件版本信息，并提供最新 PDF 格式「RedHawk Users' Manual」的访问
\end{itemize}

\subsection{主显示区}
\subsubsection*{Primary Display Area}

主显示区可显示全芯片设计或设计片段，并叠加所选特性、高亮与分析指标。

\subsection{日志显示区}
\subsubsection*{Log Display Area}

GUI 左下角的日志显示区显示 RedHawk 文本输入、进度状态、错误消息与结果。

\subsection{控制按钮}
\subsubsection*{Control Buttons}

GUI 窗口右侧的控制按钮组可改变显示或激活 RedHawk 命令。将光标置于按钮上可识别其功能。控制按钮包括：
\begin{itemize}
  \item \textbf{View} 按钮——修改主显示区视图的大小与中心
  \item \textbf{Configuration} 按钮——查看层、设置色阶范围、查看电源 pad
  \item \textbf{View Results} 按钮——在设计上叠加显示各类参数的「地图」，例如功耗密度、decap 密度、压降与电流
  \item \textbf{Query} 按钮——查询所选设计对象的属性
  \item \textbf{Coordinates} 读数——指示光标的 x,y 位置，以及主显示区边界框（单位：微米）
\end{itemize}

\subsection{设计视图区}
\subsubsection*{Design View Area}

迷你视图区（min-view area）始终显示主显示区所选内容的「Show All」视图，以展示主视图的上下文与方向。

\subsection{TCL 命令行}
\subsubsection*{TCL Command Line}

界面底部为 TCL 命令行，用于输入文本命令。语法约定与完整 TCL 命令列表见附录~\ref{chap:cmd}「TCL / Script Commands」。

\subsection{使用 GUI}
\subsubsection*{Using the GUI}

\subsubsection{导出与导入 GUI 配置}
\paragraph*{Exporting and Importing a GUI Configuration}

可通过两个 \textbf{Configuration} 按钮对话框修改的所有设置，均可保存到配置文件中供后续使用：
\begin{itemize}
  \item 菜单命令 \textbf{File $\rightarrow$ Export GUI Config}
  \item 或使用 TCL 命令：
\begin{lstlisting}
export guiconf <GUI-config-filename>
\end{lstlisting}
\end{itemize}

将 GUI 配置恢复到已保存版本有三种方式：
\begin{enumerate}
  \item 从菜单 \textbf{File $\rightarrow$ Import GUI Config} 选择要恢复的已保存配置文件
  \item 使用 TCL 命令：
\begin{lstlisting}
import guiconf <GUI-config-filename>
\end{lstlisting}
  \item 从命令行设置环境变量 \texttt{APACHE\_CONF\_FILE}：
\begin{lstlisting}
setenv APACHE_CONF_FILE <GUI-config-filename>
\end{lstlisting}
\end{enumerate}

\begin{noteBox}
方法 3 优先级最低；方法 1 与方法 2 会覆盖方法 3。
\end{noteBox}

若不想导出/导入整个 GUI 配置，可使用 TCL 命令 \cmd{config} 随时从命令行更改 viewlayers、viewnets 与 colormap。例如，\cmd{config colormap} 命令可按绝对压降或百分比配置显示的最大/最小压降，如下例：
\begin{lstlisting}
config colormap -percent -min 10 -max 24 -instance
\end{lstlisting}
这会将显示的实例最小压降设为标称 Vdd 的 10\%，最大压降百分比设为 24\%。

\section{RedHawk 数据准备（静态与动态分析）}
\subsection*{RedHawk Data Preparation - Static and Dynamic Analysis}

\subsection{RedHawk 程序文件}
\subsubsection*{RedHawk Program Files}

RedHawk 分析所需数据文件来自三类来源：
\begin{itemize}
  \item \textbf{Des}——芯片设计流程所需或生成的标准文件
  \item \textbf{Cr}——需为 RedHawk 分析自行创建的 Apache 文件
\end{itemize}

输入文件可能为静态或动态压降分析所必需，如表~\ref{tab:rh-data-files} 所示；也可能为可选——非严格必需的文件通常在提供时可提高分析精度。非必需但强烈推荐的文件标注为「Dyn-Rec」。

\begin{table}[htbp]
\centering
\caption{RedHawk 分析数据文件（表 3-1）}
\label{tab:rh-data-files}
\small
\begin{tabularx}{\textwidth}{@{}lXcc@{}}
\toprule
文件 & 说明 & 来源 & 必需 \\
\midrule
\gsr{DEF (*.def)} &
电路实例、电源/地网及其他电路元素的物理描述。当前支持的 DEF 版本为 lefdef 5.7。若有多个 \texttt{.def} 文件，必须在 GSR 关键字 \gsr{DEF_FILES}（见附录~\ref{chap:files}「DEF\_FILES」）中逐行列出各文件的绝对或相对路径（相对于 RedHawk 运行目录），并指定 \texttt{top} 或 \texttt{block} 以表明顶层或 block 级 DEF 文件。 &
Des & Stat/Dyn \\
\gsr{LEF (*.lef)} &
Library Exchange Format 文件，描述库单元的物理信息。当前支持的 LEF 版本为 lefdef 5.7。若有多个 LEF 文件，必须在 GSR 关键字 \gsr{LEF_FILES}（见附录~\ref{chap:files}「LEF\_FILES」）中逐行列出绝对或相对路径。此外，technology 文件必须包含以下四个关键字（可注释掉），RedHawk 才能将其识别为 technology LEF 文件：\texttt{LAYER}、\texttt{TYPE}、\texttt{WIDTH}、\texttt{ROUTING}。 &
Des & Stat/Dyn \\
\gsr{LIB file (*.lib)} &
Synopsys 格式库文件，包含库单元的逻辑描述、功耗与时序表，用于功耗计算。也可将全部 LIB 文件数据放在 GSR 文件的 \gsr{LIB_FILES} 关键字下。每个 operating condition 最多可处理五个字符串变量参数。 &
Des & Stat/Dyn \\
Device models &
用于仿真的 BSIM 器件模型。 & Des & Dyn \\
Standard cell SPICE netlists &
标准单元与去耦电容单元的 SPICE 网表、SPICE model card 与库 subcircuit。创建 Apache Power Library 中的电流剖面所必需。 &
Des & Dyn \\
\gsr{STA file} &
静态时序分析文件，产生实例相关的最小/最大转换时间，并定义 timing window 与时钟网数据。签核精度所必需。STA 输出文件必须在 \texttt{.gsr} 文件中用 \gsr{STA_FILE} 关键字定义。也可提高静态 IR 功耗与压降分析精度。若某实例缺少 timing window 信息，则该实例功耗假定为零，除非在 \texttt{.gsr} 文件中用 \gsr{BLOCK_POWER_FOR_SCALING} 关键字指定，或在 \gsr{VCD_FILE} 中提供。生成该文件的详细说明见第~\ref{chap:ate}~章「使用 Apache Timing Engine（ATE）创建时序文件」。 &
Des & Dyn-Rec \\
\gsr{SPEF file} &
Standard Parasitic Exchange File，实例相关的信号线寄生数据。签核精度所必需。层次化设计时，以下情况均可接受：(1) 层次 DEF 配层次 SPEF——每个 SPEF 必须关联一个 DEF block；(2) 层次 DEF 配扁平 SPEF；(3) 扁平 DEF 配扁平 SPEF。在信号完整性模式下执行 PJX 时序分析时，必须提供耦合电容数据。 &
Des & Dyn-Rec \\
\gsr{VCD file} &
Value Change Dump 文件，用于确定网翻转活动与峰值功耗。\texttt{vcdtrans}/\texttt{vcdscan} 工具用于基于 VCD 的分析所必需。若有 VCD 文件，可用 \gsr{VCD_FILE} 关键字从 VCD 文件读取网翻转信息。若无 VCD 文件，可在 \texttt{.gsr} 文件中指定默认翻转率与运行频率。图~\ref{fig:vcdscan} 为按周期功耗的 VCD 示例图。 &
Des & Dyn-Rec \\
\gsr{GDSII file} &
用于创建 DEF 视图的电路元素物理描述，由 \texttt{gds2def}/\texttt{gds2rh} 工具使用。对存储器、flip-chip bump 层描述与模拟块进行精确分析所必需——这些内容可能仅以 GDSII 格式提供，需通过 RedHawk 的 GDSII 数据准备工具 \texttt{gds2def}/\texttt{gds2rh}、\texttt{gds2def -m}/\texttt{gds2rh -m} 转换。详见附录~\ref{chap:util}「gds2rh/gds2def」。可用 GSR 关键字 \gsr{GDS_CELLS} 导入 GDS 单元，该关键字列出单元及 \texttt{gds2def}/\texttt{gds2rh} 工具的运行路径。语法见附录~\ref{chap:files}「GDS\_CELLS」。 &
Des & Dyn-Rec \\
\gsr{Technology file (*.tech)} &
设计相关文件，包含各工艺角的金属层与封装信息。technology 文件内容与格式详见附录~\ref{chap:files}。 &
Cr & Stat/Dyn \\
GDS layer map file &
实现 GDS 层号到 LEF 层名的映射。 & Cr & Stat/Dyn \\
\gsr{Global System Requirement file (*.gsr)} &
设计相关信息：工作条件、分析参数（如 VDD 网全局规格、时钟根及其频率、默认逻辑翻转率），以及复位网、I/O 网、I/O pin 输出负载等特殊网的开关条件。GSR 还包含其他输入文件的规格。确保 GSR 文件包含所有所需文件的路径与名称，例如 TECH、pad、STA、DEF、LEF、LIB、VCD、GDS、SPEF 与 package。\texttt{.gsr} 文件内容说明见附录~\ref{chap:files}。 &
Cr & Stat/Dyn \\
Pad files (\gsr{.ploc}, \gsr{.pad}, or \gsr{.pcell}) &
电源/地 pad 描述、单元名、实例名及 X,Y 位置。pad 文件说明（含 pad 单元名 \texttt{.pcell} 与 pad 实例名 \texttt{.pad} 文件）见附录~\ref{chap:files}。 &
Cr & Stat/Dyn \\
\bottomrule
\end{tabularx}
\end{table}

\begin{noteBox}
RedHawk 可读取 \texttt{*.gz}（gzip）格式的 LEF、DEF、STA、SPEF、VCD、FSDB、SAIF 与 GDS 文件。
\end{noteBox}

\figplaceholder{Figure 3-2}{vcdscan 按周期功耗曲线}
\label{fig:vcdscan}

\subsection{多路 Vdd/Vss 分析}
\subsubsection*{Multiple Vdd/Vss Analysis}

RedHawk 原生支持多电源/地域，不仅能准确处理具有多路 Vdd 供电的设计，还能处理具有多路电源/地 pin 的单个单元，并始终查找支持每单元多域的数据。多 Vdd 设计元素的常见示例如下：
\begin{itemize}
  \item 电平转换器（level shifters）
  \item 具有主电源与 retention 电源的触发器与锁存器
  \item 用于将不同电位 block 隔离的隔离逻辑
  \item 阵列与外围部分具有独立电压的存储器/寄存器文件 block
  \item 同时具有外部与内部电源/地供电的 I/O 单元
\end{itemize}

确定多电源域单元在静态与动态分析中正确建模所需的数据，见以下各段说明。

对于具有单路 VDD 与 VSS 供电 pin 的单元，RedHawk 配置文件可能需要修改，但输入数据无需更改。在大多数情况下，仅使用单电源/地域单元的设，只需重建任何经 \texttt{gds2def -m}/\texttt{gds2rh -m} 建模的单元。这并非多 Vdd 支持的结果，而是 \texttt{gds2def -m}/\texttt{gds2rh -m} 建模 IP 中针对非均匀电流分布的更新建模结构所致。

对于具有多路 VDD pin 以及一路或多路 VSS pin 的单元，RedHawk 需要额外的用户数据与配置文件。

\subsubsection{输入数据文件差异摘要}
\paragraph*{Summary of Differences in Input Data Files}

表~\ref{tab:mvdd-input} 汇总多 Vdd 分析所需数据文件变更。标注「GEN」的条目与 RedHawk 数据处理的一般变更相关，与多 Vdd 无直接关系。

\begin{table}[htbp]
\centering
\caption{多 Vdd 分析所需输入数据变更（表 3-2）}
\label{tab:mvdd-input}
\small
\begin{tabularx}{\textwidth}{@{}llX@{}}
\toprule
文件类型 & 需要变更的单元类型 & 所需变更（GEN/MVdd） \\
\midrule
LIB & 仅多域 &
MVdd：库、单元与 pin 级定义中需要特定属性。Liberty 格式不接受具有多个 ground pin 的单元，此时必须创建 custom LIB 文件。 \\
LEF & 多域；单域 &
GEN：所有电源与地 pin 必须正确定义 \texttt{USE < >} 与 \texttt{DIRECTION < >} 属性。 \\
User instance power file & 仅多域 & MVdd：按 pin 指定功耗。 \\
APL data & 仅多域 &
MVdd：重新运行 APL 表征，为每个 Vdd 与 Vss arc 创建独立的实例电流剖面，以及每 pin 的 Cdev。 \\
Memory models & 多域；单域 &
GEN：使用 \texttt{gds2def -m}/\texttt{gds2rh -m} 创建的存储器模型应重新提取，以使用增强的存储器处理能力。若使用 \texttt{gds2def}/\texttt{gds2rh}，则无需更改。 \\
\bottomrule
\end{tabularx}
\end{table}

表~\ref{tab:mvdd-config} 汇总执行多 Vdd/Vss 域分析所需的 RedHawk 配置文件变更。表后各段进一步说明所需变更。

\begin{table}[htbp]
\centering
\caption{多 Vdd 配置文件变更（表 3-3）}
\label{tab:mvdd-config}
\small
\begin{tabularx}{\textwidth}{@{}llX@{}}
\toprule
文件类型 & 需要变更的单元类型 & 所需变更（GEN/MVdd） \\
\midrule
GSR & 多域 &
GEN：使用基于 GSR 的规格指定所有输入文件。 \\
GSR & 单域 &
使用 \gsr{GDS_CELLS} 关键字，在使用 GDS2RH -m 时采用新存储器模型。 \\
Custom LIB & 多域且多 GND pin & MVdd：指定电源/地 arc 对。 \\
APL configuration & 多域 & MVdd：为 APL-DI 指定 LEF 文件。 \\
AVM configuration & 多域；单域 &
MVdd：在配置文件中指定 VDD 与 GND pin 名称，如下：\gsr{VDD_PIN <net names>}、\gsr{GND_PIN <net names>}。 \\
GSC & 多域；单域 &
GEN：\texttt{OFF} 状态仅用于电源门控 block。要将普通单元置为「Off」，请使用 \texttt{DISABLE}。 \\
\bottomrule
\end{tabularx}
\end{table}

\subsubsection{多 Vdd/Vss 域设计的 Liberty 库语法差异}
\paragraph*{Liberty Library Syntax Differences for Multiple Vdd/Vss Domains}

以下各节说明多 Vdd/Vss 域设计对 LIB 文件的要求。

\subparagraph{\texttt{power\_supply} 组（库级）}
\texttt{power\_supply} 组捕获所有关于电压变化的标称信息，定义默认电源名与各电源的标称电压值。新语法示例如下：

语法：
\begin{lstlisting}
power_supply () {
   default_power_rail : <domain_name> ;
   power_rail (<domain_name>, <voltage>) ;
   ...
}
\end{lstlisting}

示例：
\begin{lstlisting}
power_supply () {
    default_power_rail : VDD ;
    power_rail(VDDNW,0.81);
    power_rail(VDDC,0.81);
    power_rail(VDD,0.81);
    power_rail(VDDIN,0.95);
}
\end{lstlisting}

\subparagraph{\texttt{rail\_connection} 语句（单元段）}
\texttt{rail\_connection} 属性必须定义单元中使用的所有多路电源与地域 pin，将库电源域名映射到其所连接的单元 LEF 电源 pin 名。新语法示例如下：

语法：
\begin{lstlisting}
rail_connection ( <power_pin_name>, <power_domain_name>) ;
rail_connection ( <grnd_pin_name>, <grnd_domain_name> ) ;
...
\end{lstlisting}

示例：
\begin{lstlisting}
rail_connection (VDDS, VDD1) ;
rail_connection (VDD_RET, VDD2);
rail_connection (VSS_A, VSS) ;
\end{lstlisting}

\subparagraph{\texttt{leakage\_power} 语句}
\texttt{leakage\_power} 语句为内部/漏电功耗定义 rail 相关的功耗值。新语法示例如下：
\begin{lstlisting}
leakage_power () {
      value : 2.500000E+00;
      power_level : "VDDNW";
    }
leakage_power () {
      value : 6.740750E+03;
      power_level : "VDD";
    }
\end{lstlisting}

\subparagraph{\texttt{input\_signal\_level} 属性}
\texttt{input\_signal\_level} 属性定义应施加到输入或 inout pin/bus/bundle 的电压电平。

语法：\texttt{input\_signal\_level : <domain\_name> ;}

\subparagraph{\texttt{output\_signal\_level} 属性}
\texttt{output\_signal\_level} 属性定义应施加到 inout 或输出 pin/bus/bundle 的电压电平。

语法：\texttt{output\_signal\_level : <domain\_name> ;}

\subparagraph{\texttt{internal\_power}}
\texttt{internal\_power} 属性表为所有电源 arc 定义开关期间的短路能量值。

\subparagraph{\texttt{power\_level}}
\texttt{power\_level} 属性指定用于表征 pin 内部功耗的电源。

语法：\texttt{power\_level : "<domain\_name>" ;}

示例：
\begin{lstlisting}
pin ( Z ) {
    direction : input ;
    output_signal_level : VDD1;
    internal_power() {
         power_level : VDD1 ;
         power (power_template) {
           values ("1, 2, 3, 4");
         }
        }
    } /* end pin */
\end{lstlisting}

\subsubsection{Custom LIB 文件中的 P/G Arc 定义}
\paragraph*{P/G Arc Definitions in Custom LIB Files}

为支持具有多路 Vdd 与多路 Vss pin 的单元，必须指定 P/G arc，定义每对 VDD 节点到相关 GND 节点之间的电流路径。当存在多电源域与多地域时，定义电源域 arc 可向静态与动态仿真器提供正确分配电流所需的电流节点对。

若有 UPF LIB，可从 LIB 中的 \texttt{related\_power\_pin} 与 \texttt{related\_ground\_pin} 值获取 P/G arc 信息。否则，对于多 Vdd/Vss 单元，必须在 custom LIB 中提供 P/G arc 数据。

\texttt{pgarc} 对象可在设计级 custom LIB 中定义如下：
\begin{lstlisting}
pgarc {
      <vdd_pin_name> <vss_pin_name>
      ...
}
\end{lstlisting}
此时定义适用于设计中所有多电源供电单元。\texttt{pgarc} 也可在单元内定义，仅适用于该单元：
\begin{lstlisting}
cell <cell_name> {
pgarc {
      <vdd_pin_name> <vss_pin_name>
     ...
        }
}
\end{lstlisting}

\textbf{custom library 中通常不需要 P/G arc 定义：}
\begin{itemize}
  \item 1 个 power 与 N 个 ground pin（$N = 1$ 或更多）。若一路 power 连接到多个 ground，则假定 power pin 的电流在各 ground arc 间均分。注意 APL 捕获的是 VDD 电流，而非 VSS 电流。
  \item N 个 power 与 1 个 ground，且所有 N 个 power pin 在电流文件（或 LIB 功耗）中均有 profile。
\end{itemize}

\textbf{custom library 中需要 P/G arc 定义：}
\begin{itemize}
  \item N 个 power 与 M 个 ground pin（$N > 1$，$M > 1$）
  \item N 个 power 与 1 个 ground pin——在存储器 IP 中很常见——其中部分 power pin 在电流文件（或 LIB 中缺少功耗）中没有 profile，以便现实地分配 VDD 电流
\end{itemize}

例如，假设库中大多数单元具有三个电源域 VDD、VDD2 与 VDDL，以及一个地域 VSS。但单元 \texttt{ram\_mvdd} 具有三个电源域 VDD、VDD2 与 VDDL，以及两个地域 VSS 与 VSS2。在 \texttt{ram\_mvdd} 单元中，VDD 与 VDDL 均使用 VSS 作为地，而 VDD2 使用 VSS2 作为地。由于仅有一个地域，主设计的 pgarc 关系不需要 custom LIB 文件；但 \texttt{ram\_mvdd} 单元同时具有两个地域与三个电源域，必须在 custom LIB 文件中定义 pgarc 关系。custom 文件格式示例如下：
\begin{lstlisting}
pgarc {
     VDD VSS
     VDD2 VSS
     VDDL VSS
}
cell ram_mvdd {
  pgarc {
     VDD VSS
     VDD2 VSS2
     VDDL VSS
     }
}
\end{lstlisting}

\begin{noteBox}
若某单元的 p/g arc 在 custom lib 中未完整定义，RedHawk 会忽略该单元在 LEF/LIB 中的其余 p/g arc。若 custom lib 中未为该单元定义 arc，RedHawk 可从 LEF/LIB 隐式读取该信息；但隐式 p/g arc 可能因 LEF 中的 pin 类型或 LIB 中不完整的 p/g arc 信息而不正确。因此必须确保 custom lib 中定义的 p/g arc 完整且正确——它们具有最高优先级，会覆盖 LEF/LIB 中的 p/g arc 信息。此外，custom lib 还可灵活地将电流仅分配到您关心的电源/地 pin。在 GSR 文件中，\gsr{LIB_FILES} 关键字用于指定 custom LIB 文件名（仅可指定一个 custom \texttt{*.lib} 文件）：
\begin{lstlisting}
LIB_FILES {
<lib_filename> CUSTOM
...
}
\end{lstlisting}
若不确定哪些单元需要在 custom library 中指定，可运行 \cmd{setup design}，然后查看 \texttt{adsRpt/apache.refCell.noPGArc} 文件中的数据。RedHawk 检测具有多个 ground pin 且无 custom LIB 文件的单元，并在报告 \texttt{adsRpt/apache.refCell.noPGArc} 中列出，同时给出每个单元的所有电源与地 pin。

示例：
\begin{lstlisting}
#<cell_name> <vdd_pin_names> <gnd_pin_names>
<SC_ANALOG> <VDD1A VDD2A> <VSS1A VSS2A>
\end{lstlisting}
\end{noteBox}

\subparagraph{状态相关漏电功耗（state dependent leakage power）}
RedHawk 可输入 custom lib 文件，详细说明哪些状态/布尔组合应包含/排除在（状态相关）漏电功耗计算中。

例如，某些状态下状态相关漏电功耗可能非常高，尤其是存储器。该状态通常是不在常规模式中使用的状态。可通过 custom lib 设置，不将该状态的漏电功耗计入。

语法：
\begin{lstlisting}
cell <cell_name> {
       exclude_states {
           pin <active_pin_name> {
                <boolean expression1>
                <boolean expression2>
                ..
           }
           leakage_power {
                <boolean expression1>
                <boolean expression2>
                ..
           }
     }
     include_states {
         pin <active_pin_name> {
              <boolean expression1>
              <boolean expression2>
              ..
         }
         leakage_power {
              <boolean expression1>
              <boolean expression2>
              ..
         }
     }
}
\end{lstlisting}

单元名定义与布尔表达式中接受通配符。

示例：
\begin{lstlisting}
cell OR3* {
 exclude_states {
 leakage_power {
 "*!C"
      }
            }
 include_states {
 leakage_power {
 "*&C"
      }
            }
                  }
\end{lstlisting}

\subsubsection{多 Vdd 域设计的 GSR 关键字}
\paragraph*{GSR Keywords for Multi-Vdd Domain Designs}

\subparagraph{P/G arc 错误}
可使用关键字 \gsr{IGNORE_PGARC_ERROR 1} 抑制 P/G arc 错误。

\subparagraph{\gsr{BLOCK_POWER_FOR_SCALING} 与 \gsr{BLOCK_POWER_FOR_SCALING_FILE} 关键字}
以下关键字语法应用于在 \gsr{BLOCK_POWER_FOR_SCALING} 定义文件中处理具有单路或多路 Vdd/Vss 供电的单元。但在 6.1 及更早版本中，\texttt{fullchip <full\_mVdd\_inst\_path>} 选项不能用于 GSR 的 \gsr{BLOCK_POWER_FOR_SCALING} 条目。该关键字用法见附录~\ref{chap:files}「BLOCK\_POWER\_FOR\_SCALING」。

\gsr{BLOCK_POWER_FOR_SCALING} 文件语法：
\begin{lstlisting}
{ CELLTYPE <cell_name> <cell_power_Watts> ? <Vdd_pin>?
                 ? <cell_current_A> <Vss_pin> ?
...
[FULLCHIP |<top_block_name>] <full_block_instance_path> <pwr_Watts>
...
 fullchip <full_mVdd_inst_path> <power_Watts> ?<VDD_pin>?
                 ? <current_A> <Vss_pin> ?
...
}
\end{lstlisting}

\subsubsection{多 Vdd 设计的 APL 要求}
\paragraph*{APL Requirements for Multi-Vdd Designs}

多 Vdd/Vss 域单元 APL 表征的关键要素：
\begin{itemize}
  \item Design Dependent 表征无需更改 APL 配置文件。若提供给 RedHawk 的 LEF/LIB 文件完整，RedHawk 会自动生成 mVdd 表征所需文件。
  \item APL-DI 在重新表征时需要新关键字 \gsr{LEF_FILES}。对于 APL-DI，必须在 APL 配置文件中列出覆盖所有待表征多 Vdd 单元的 LEF 文件，语法为：
\begin{lstlisting}
LEF_FILES {<fileA> ... <fileN> }
\end{lstlisting}
  \item RedHawk 使用 LEF 文件推导 P/G arc 信息。为每个单元中定义的电源/地 pin arc 计算电流剖面与 decap 值。
\end{itemize}

\begin{noteBox}
\texttt{*.current} 与 \texttt{*.cdev} 文件为二进制格式。支持导入并合并混合 RedHawk 5.x 版本的 APL 文件，但有一定限制。GSR 文件中的 \gsr{APL_FILES} 关键字结构为：
\begin{lstlisting}
APL_FILES {
<APL_binary_filename>            current
<APL_binary_filename>            cdev
<Avm.conf_filename>              avm
<AVM_binary_filename> current_avm
<AVM_binary_filename> cap_avm
<APL binary filename> pwcap
...
}
\end{lstlisting}
在 7.x 版本中，来自 5.x 与 6.x 的 \texttt{<cell>.current} 文件可在 APL 导入与合并时混合使用。更多信息见第~\ref{chap:apl}~章「导入 APL 文件」。
\end{noteBox}

\subsubsection{AVM 配置文件}
\paragraph*{AVM Configuration File}

使用附录~\ref{chap:apl}「LIB2AVM Configuration File」中详述的语法，为单路或多路 Vdd/Vss 设计指定 power/gnd pin。

\subsection{数据准备与自动化 \texttt{rh\_setup.pl} 脚本}
\subsubsection*{Data Prep and Using the Automated \texttt{rh\_setup.pl} Script}

本节说明如何指定必要的输入设计文件、创建配置与运行时命令文件以运行 RedHawk（使用自动化 setup 脚本），并描述输入与工作文件的推荐目录结构。

\subsubsection{推荐的 RedHawk 目录结构}
\paragraph*{Recommended RedHawk Directory Structure}

每个设计建议采用以下 RedHawk 目录结构：

\begin{verbatim}
<design_name>/
  design_data/          # 所有用户提供的输入数据
    def/                # 所有 *.def（*.gz 亦可）
    lef/                # 所有 *.lef
    lib/                # 所有 *.lib
    tech/               # RedHawk technology 文件
    ploc/               # 电压位置或 pad cell 文件
    timing/             # 时序数据（动态分析需要，可选）
    spf/                # 全芯片或 block 级 SPEF/DSPF
    vcd/                # VCD 文件
    power/              # 实例级功耗数据
    gds/                # 设计中所有存储器/宏的 GDS 文件
    gds2def/ 或 gds2rh/ # GDS 转换结果文件
    apl/                # APL/表征数据目录
    memory/             # HSIM/Spice 波形
    avm/                # AVM 结果
  <RH_run1>/            # 工作目录
    static/             # 静态分析结果
    setupApl/           # 创建 APL 数据文件
    dynamic-dvd/        # vectorless 方法动态分析结果
    dynamic-vcd/        # VCD 输入方法动态分析结果
    adsPower/           # RedHawk 创建的功耗计算数据
    adsRpt/             # Cmd、Error、Warn、Log 文件（分文件夹）
    .apache/            # RedHawk 运行工作文件
  <RH_run2>/
  <RH_runN>/
\end{verbatim}

各目录内容说明：
\begin{itemize}
  \item \texttt{design\_data/}：所有用户提供的输入数据目录
  \begin{itemize}
    \item \texttt{def/}：设计中所有 DEF 文件，扩展名为 \texttt{*.def}（\texttt{*.gz} 压缩文件亦可，RedHawk 可处理）
    \item \texttt{lef/}：所有 LEF 文件，扩展名为 \texttt{*.lef}
    \item \texttt{lib/}：所有 LIB 文件，扩展名为 \texttt{*.lib}
    \item \texttt{tech/}：RedHawk technology 文件
    \item \texttt{ploc/}：电压位置或 pad cell 文件
    \item 以下 \texttt{design\_data/} 子目录为可选：\texttt{timing/}（动态分析所需时序数据）、\texttt{spf/}（全芯片或 block 级 SPEF/DSPF）、\texttt{vcd/}、\texttt{power/}、\texttt{gds/}、\texttt{gds2def/} 或 \texttt{gds2rh/}、\texttt{apl/}、\texttt{memory/}、\texttt{avm/}
  \end{itemize}
  \item \texttt{<RH\_runX>/}：工作目录
  \begin{itemize}
    \item \texttt{static/}：静态分析结果目录
    \item \texttt{setupApl/}：创建 APL 数据文件的目录
    \item \texttt{dynamic-dvd/}：vectorless 方法动态分析结果
    \item \texttt{dynamic-vcd/}：VCD 输入方法动态分析结果
    \item \texttt{adsPower/}：RedHawk 创建的功耗计算数据
    \item \texttt{adsRpt/}：RedHawk 创建的 Cmd、Error、Warn、Log 文件（分文件夹存放），以及各类型最新文件的链接
    \item \texttt{.apache/}：RedHawk 运行工作文件
  \end{itemize}
\end{itemize}

\begin{noteBox}
在同一运行目录中开始新运行前，\texttt{.apache} 目录中任何先前 RedHawk 运行的工作文件都会被删除。
\end{noteBox}

工作目录 \texttt{static/}、\texttt{dynamic-dvd/}、\texttt{dynamic-vcd/} 可自定名称，但应与 \texttt{design\_data/} 目录处于同一层级。将所需输入文件复制到 \texttt{design\_data/} 目录，为设置与运行 RedHawk 做准备。

\begin{noteBox}
若主目录或当前工作目录中存在名为 \texttt{.redhawkrc} 的文件，且其中包含一条或多条标准 RedHawk TCL 命令，则 RedHawk 会在执行任何其他 RedHawk 命令之前先执行该文件。
\end{noteBox}

为获得自动化 setup 脚本的最佳运行效果，\texttt{design\_data/} 目录应采用上述结构，但并非强制。只要指定了路径，setup 脚本即可找到所需文件。

\subsubsection{自动化脚本 \texttt{rh\_setup.pl}}
\paragraph*{Automated Script rh\_setup.pl}

\subparagraph{设置自动化 TCL Setup 脚本}
要运行 RedHawk，必须准备适当的 Global System Requirements（GSR）文件，并编写运行各所需步骤的命令脚本。为自动化这些任务，RedHawk 提供 \texttt{rh\_setup.pl} 工具，协助检查所需文件、设置 GSR 文件，并准备首次运行 RedHawk 的命令文件。使用 setup 脚本的主要特性：
\begin{itemize}
  \item \texttt{rh\_setup.pl} 创建启动 RedHawk 仿真（静态、动态或特殊评估）所需的全部文件：
  \begin{itemize}
    \item \texttt{<design>.gsr}——Global System Requirements 关键字设置
    \item \texttt{run\_static.tcl}——静态运行脚本
    \item \texttt{run\_dynamic.tcl}——动态运行脚本
    \item \texttt{run\_signalEM.tcl}——信号 EM 分析运行脚本
    \item \texttt{run\_lowpower.tcl}——低功耗分析 ramp-up 运行脚本
  \end{itemize}
  \item 生成 GSR 文件所用的参数值可来自四个来源，按优先级从高到低：
  \begin{enumerate}
    \item \texttt{rh\_setup.pl} 命令行参数
    \item 保存在 \texttt{rh\_setup.init} 文件中的先前 \texttt{rh\_setup.pl} 调用参数
    \item 由变量 \texttt{\$APACHEDA\_TEMPLATE\_GSR} 指向的模板 GSR 文件（若存在）
    \item RedHawk 预设的默认值
  \end{enumerate}
  \item 无需记忆 GSR 关键字的精确语法——\texttt{rh\_setup} 生成的 GSR 关键字语法正确
  \item 若采用推荐的数据目录结构，脚本会自动定位文件。例如，若在 \texttt{../../data/def/*.def} 目录中找到 DEF 文件，则全部写入生成的 GSR
  \item 脚本支持 UNIX 通配符在非标准位置搜索所需文件。例如，使用文件搜索模式 \texttt{-def *.def */*.def */*/*.def}，可找到最多两层子目录下的所有 DEF 文件。若同一文件名出现在多个目录中，则使用首先指定目录中的文件
  \item 中央 CAD 工程组可通过「模板 GSR」文件为特定 RedHawk 参数（例如所需动态仿真步长）提供默认值。运行 setup 脚本前，使用全局环境变量 \texttt{\$APACHEDA\_TEMPLATE\_GSR} 指向模板 GSR 文件
  \item 脚本输入可增量进行。可用一个或多个命令行参数调用脚本；若缺少任何必需参数（例如电源网名），脚本会提示。可重新调用脚本并添加或更改参数，直至所有必需信息完整。然后 \texttt{rh\_setup.pl} 检查所有输入文件是否存在，并创建上述文件列表
\end{itemize}

\subparagraph{使用 rh\_setup 工具}
使用 setup 脚本的步骤如下：

\begin{enumerate}
  \item 脚本 \texttt{rh\_setup.pl} 的 UNIX 命令行格式为：
\begin{lstlisting}
rh_setup.pl <command_options> [<values>]
...
\end{lstlisting}
  \item 不带选项执行脚本命令：
\begin{lstlisting}
rh_setup.pl
\end{lstlisting}
  将显示必需选项、参数值与文件的定义，以及可选项。这些选项为：
  \begin{itemize}
    \item \texttt{-top\_cell}（必需）——设计中顶层单元名，与对应 DEF 文件中指定的一致（见 DEF 文件的 \texttt{DESIGN} 段）
    \item \texttt{-mode}——分析模式，\texttt{early\_analysis} 或 \texttt{sign\_off\_analysis}。默认 \texttt{sign\_off\_analysis}
    \item \texttt{-analysis}——分析类型：\texttt{static}、\texttt{dynamic}、\texttt{low\_power} 或 \texttt{signalEM}。默认创建 static 与 dynamic 的 TCL 文件
    \item \texttt{-vdd\_nets}（必需）——指定 Vdd 网名与标称电压值
    \item \texttt{-vss\_nets}——Vss 网名（不指定电压）。静态或 signalEM 运行可选；动态或 lowpower 运行必需
    \item \texttt{-frequency}（必需）——定义设计的主导工作频率（Hz），即消耗设计中大部分功耗的频率。若有多个主导频率，指定较低频率。见附录~\ref{chap:files}「FREQUENCY」
    \item \texttt{-tech\_file}（必需）——定义 \texttt{tech/} 目录中的介质、金属层参数与 via。若未指定，则搜索 \texttt{rh\_setup.defaults} 中的默认路径
    \item \texttt{-def\_files}——若所需文件在 \texttt{def/} 目录中则可选。指定顶层模块、层次与 IP block 的 DEF 文件
    \item \texttt{-lef\_files}——若所需文件在 \texttt{lef/} 目录中则可选。指定工艺、库单元、层次与 IP block 的文件
    \item \texttt{-lib\_files}——若所需文件在 \texttt{lib/} 目录中则可选。指定 Synopsys Liberty 格式（\texttt{.lib}）时序库。早期分析或静态运行不需要；动态、lowpower 与签核运行需要
    \item \texttt{-pad\_files}（必需）——自动化流程从 \texttt{ploc/} 目录获取文件，或指定 \texttt{def\_pins} 选项使用 DEF pin 作为 block 级分析的电源 pad（强制 RedHawk 在顶层 DEF 文件中描述的每个 Vdd/Vss pin 中心放置一个 pad）。RedHawk PLOC 文件提供理想电压源的 x,y 坐标。需要其中一种 pad 文件
    \item \texttt{-sta\_files}——若所需文件在 \texttt{timing/} 目录中则可选。指定来自 STA 的时序文件，见第~\ref{chap:ate}~章。签核分析运行需要
    \item \texttt{-spf\_files}——若所需文件在 \texttt{spf/} 目录中则可选。指定布局后寄生的层次 SPEF 或 DSPF 文件
    \item \texttt{-apl\_files}——包含 APL 表征电流剖面的 RedHawk \texttt{cell.current} 文件。模式为 \texttt{sign\_off} 且分析类型为 \texttt{dynamic} 或 \texttt{lowpower} 时必需
    \item \texttt{-aplcap\_files}——来自 APL 表征的 RedHawk \texttt{cell.cdev} 器件电容文件。模式为 \texttt{sign\_off} 且分析类型为 \texttt{dynamic} 或 \texttt{lowpower} 时必需
    \item \texttt{-aplpwcap\_files}——来自 APL 上电表征的 RedHawk \texttt{cell.pwcap} PWL decap 文件。模式为 \texttt{sign\_off} 且分析类型为 \texttt{lowpower} 时必需
    \item \texttt{-gds\_dirs}——若所需文件在 \texttt{gds2def/}、\texttt{gds2rh/} 目录中则可选。指定 \texttt{gds2def} 或 \texttt{gds2rh} 在指定目录中创建的所有 LEF、DEF 与 PRATIO 文件
    \item \texttt{-toggle}（可选）——自动化流程分配默认值
    \item \texttt{-input\_transition}（可选）——自动化流程分配默认值
  \end{itemize}

  \begin{noteBox}
  可缩写命令选项，只要无歧义即可。例如，可用 \texttt{-top} 作为 \texttt{-top\_cell} 的缩写。
  \end{noteBox}

  \item 在 \texttt{design\_data/} 目录中设置好输入数据后，可在任意位置运行 \texttt{rh\_setup} 命令，指定所需值，例如：
\begin{lstlisting}
rh_setup.pl -top_cell <name of the design>
    -vdd_nets <name of VDD net> <value>
    -vss_nets <names of VSS nets>
    -pad_files <name of pad/pcell file in ploc/ dir,
    OR specify def_pins for DEF pin-based analysis>
    -def_files <path and filenames>
    -lef_files <path and filenames>
    -lib_files <path and filenames>
    -tech_file <path and filenames>
\end{lstlisting}
  假设找到所有输入信息，此次调用将创建必要的运行文件。

  \item \texttt{rh\_setup.pl} 工具将所有输入值保存在名为 \texttt{rh\_setup.init} 的文本文件中（保存在调用该工具的同一目录），因此可多次调用该工具，每次提供一个或多个必需参数，直至满足所有数据要求。此时才创建运行文件。也可直接编辑 \texttt{rh\_setup.init} 文本文件。

  \item 要执行 \texttt{rh\_setup.pl} 生成的命令脚本，使用 TCL 的 \texttt{-f} 选项，例如：
\begin{lstlisting}
redhawk -f [ run_static.tcl | run_dynamic.tcl ]
\end{lstlisting}
\end{enumerate}

\subparagraph{特殊 GSR 关键字}
部分重要 GSR 关键字维护默认值：
\begin{itemize}
  \item \gsr{DYNAMIC_SIMULATION_TIME}：\texttt{2.56e9}
  \item \gsr{DYNAMIC_TIME_STEP}：\texttt{early\_analysis} 运行为 50\,ps；\texttt{sign\_off} 运行为 20\,ps
  \item \gsr{INPUT_TRANSITION}：200\,ps
  \item \gsr{TOGGLE_RATE}：0.3（非时钟）
\end{itemize}

以下 GSR 关键字仅按分析模式生效，或在 GSR 文件中保持注释状态：
\begin{itemize}
  \item \gsr{BLOCK_POWER_ASSIGNMENT}——在原型分析期间开启，并提示用户编辑其需求
  \item \gsr{SWITCH_MODEL_FILE}——低功耗分析运行期间由 \texttt{aplsw} 工具自动生成并添加到 GSR 段
\end{itemize}

\subsection{使用 \texttt{gds\_setup.pl} 脚本进行 GDS2DEF 设置}
\subsubsection*{GDS2DEF Setup with the \texttt{gds\_setup.pl} Script}

RedHawk \texttt{gds\_setup.pl} 脚本有助于自动化生成必要配置文件，并以串行或并行方式对多个单元运行 \texttt{gds2def} 提取工具。脚本接受来自两个来源的用户数据：
\begin{itemize}
  \item \texttt{\$cwd/gds\_setup.init} 文件
  \item 命令行参数（脚本会将其保存到 \texttt{gds\_setup.init}）
\end{itemize}

与 \texttt{rh\_setup} 一样，脚本输入可增量进行。首先，脚本检查 LEF 与 GDS 文件，确保待提取的单元同时具有 GDS 与 LEF 描述。可用一个或多个选项调用脚本，然后随时重新调用以添加或更改信息，直至所有必需的 \texttt{gds2def} 设置信息完整。

\subsubsection{使用 gds\_setup 工具}
\paragraph*{Using the gds\_setup Utility}

使用 setup 脚本的步骤如下：

\begin{enumerate}
  \item 脚本 \texttt{gds\_setup.pl} 的 UNIX 命令行格式为：
\begin{lstlisting}
gds_setup.pl <command_options> [<values>]
...
\end{lstlisting}
  \item 不带选项执行脚本命令：
\begin{lstlisting}
gds_setup.pl
\end{lstlisting}
  将显示必需选项、参数值与文件的定义，以及可选项。这些选项为：
  \begin{itemize}
    \item \texttt{-vdd\_nets}（必需）——指定 Vdd 网名。
    
    示例：\texttt{-vdd\_nets VDD vdda vddb, VDD\_SOC}

    GDS 中标记为 \texttt{vdda} 与 \texttt{vddb} 的段在生成的 DEF 文件中显示为 \texttt{VDD}。\texttt{VDD\_SOC} 不重命名。
    \item \texttt{-vss\_nets}（必需）——指定 Vss 网名。
    
    示例：\texttt{-vss\_nets VSS vssa vssb, VSS\_SOC}

    GDS 中标记为 \texttt{vssa} 与 \texttt{vssb} 的段在生成的 DEF 文件中显示为 \texttt{VSS}。\texttt{VSS\_SOC} 不重命名。
    \item \texttt{-cell\_names}——指定要处理的单元名。未指定时，\texttt{gds2def} 对同时提供 LEF 与 GDS 数据的单元运行。
    \item \texttt{-map\_file}（必需）——指定 layer map 文件位置。
    \item \texttt{-lef\_files}（必需）——指定各个 LEF 文件，或指定 \texttt{*.lef} 以纳入所有 LEF 文件。
    
    示例：\texttt{-lef\_files ../lef/hdsg1\_10240x32cm16.lef ../lef/hdsg1\_2048x40cm8bw.lef}

    或：\texttt{-lef\_files ../*.lef}
    \item \texttt{-gds\_files}（必需）——指定各个 GDS 文件，或指定 \texttt{*.gds} 以纳入所有 GDS 文件。
    
    示例：\texttt{-gds\_files ../gds/hdsg1\_10240x32cm16.gds ../gds/hdsg1\_2048x40cm8bw.gds}

    或：\texttt{-gds\_files ../gds/*.gds}
    \item \texttt{-mem [ on | off ]}——设为 \texttt{on} 表示 \texttt{gds2def} 运行针对存储器。默认 \texttt{off}。
    \item \texttt{-start\_layer}——指定提取的起始层。
    
    示例：\texttt{-start\_layer m1}
    \item \texttt{-core\_start\_layer}——指定 core 提取的起始层。
    
    示例：\texttt{-core\_start\_layer m3}
    \item \texttt{-options\_file}——指定包含额外 \texttt{gds2def} 选项列表的输入文件。
    \item \texttt{-norun [ on | off ]}——设为 \texttt{on} 时仅创建配置文件，不运行 \texttt{gds2def}。默认 \texttt{off}。
    \item \texttt{-bsub\_run [ on | off ]}——设为 \texttt{on} 时，将作业提交到 LSF 集群，在多台机器上并行执行。默认 \texttt{off}。
    \item \texttt{-bsub\_command\_file}——指定包含 \texttt{bsub} 命令及定义多机执行要求的额外选项的文件。若未指定，则仅使用 \texttt{bsub} 作为命令。
  \end{itemize}
  \item 按需要编辑并更新 setup 文件选项值（见图~\ref{fig:gds-setup}），然后运行 \texttt{gds2def}。
\end{enumerate}

示例调用：
\begin{lstlisting}
gds_setup.pl -vdd_nets VDD vdd vddx vddy, VDD_SOC \
              -vss_nets VSS gnd vss \
              -map_file gds_layer.map \
              -cell_names MEM22x64 MEM22x32 \
              -lef_files lef_dir/*.lef \
              -gds_files gds_dir/*.gds \
              -mem on \
              -norun on \
              -bsub_run on \
              -bsub_command_file bsub_options.txt
\end{lstlisting}

\figplaceholder{Figure 3-3}{gds\_setup.pl 示例文件}
\label{fig:gds-setup}

\subparagraph{输出}
\texttt{OUTPUT} 目录包含 \texttt{gds2def} 运行生成的 DEF、LEF 与 PRATIO 文件；\texttt{LOG} 目录包含单元专用日志目录，其中又有日志文件。

\texttt{top\_report} 文件包含各单元是否通过 \texttt{gds2def} 处理的信息，示例如图~\ref{fig:gds-top-report}：
\begin{verbatim}
-------------------------------------------------------------------
         hdsg1_10240x32cm16 Done
         rfsg1_256x32cm4                     Done
         rfsg2_44x128cm2bw                   Fail
         ...
-------------------------------------------------------------------
\end{verbatim}

\figplaceholder{Figure 3-4}{gds2def 结果 top\_report 示例}
\label{fig:gds-top-report}

\subsection{手动导入设计数据}
\subsubsection*{Manually Importing Design Data}

上一节说明如何使用脚本准备运行 RedHawk 的输入文件。若不想使用脚本，可使用 TCL 命令行或 GUI 手动准备设计数据与配置文件。以下各节描述这些命令与流程。

\subsubsection{命令行数据导入}
\paragraph*{Command Line Data Import}

加载设计数据的 TCL 命令见表~\ref{tab:import-cmds}。注意：将所有设计数据放入 GSR 文件后导入 GSR，是执行全部设计数据输入的推荐方法。GSR 文件内容与语法详见附录~\ref{chap:files}「Global System Requirements File (\texttt{*.gsr})」。

\begin{table}[htbp]
\centering
\caption{加载设计数据的 TCL 命令（表 3-4）}
\label{tab:import-cmds}
\small
\begin{tabularx}{\textwidth}{@{}lX@{}}
\toprule
TCL 命令 & 用途 \\
\midrule
\cmd{import gsr <InputFileName>} &
导入 \texttt{.gsr} 文件中的设计信息。 \\
\cmd{import tech <InputFileName>} &
导入 \texttt{.tech} 文件。（建议在 GSR 文件中使用 \gsr{TECH_FILE} 关键字。） \\
\cmd{import lef <InputFileName>} &
导入 \texttt{.lef} 文件。（建议在 GSR 文件中使用 \gsr{LEF_FILES} 关键字。） \\
\cmd{import lib <InputFileName>} &
导入 \texttt{.libs} 文件。（建议在 GSR 文件中使用 \gsr{LIB_FILES} 关键字。） \\
\cmd{import def <InputFileName>} &
导入 \texttt{.def} 文件。（建议在 GSR 文件中使用 \gsr{DEF_FILES} 关键字。） \\
\cmd{import pad <InputFileName>} &
导入 pad 单元 \texttt{.pcell}、pad 实例 \texttt{.pad} 或 pad 位置 \texttt{.ploc} 文件。（建议在 GSR 文件中使用 \gsr{PAD_FILES} 关键字。） \\
\cmd{import guiconf <InputFileName>} &
导入 GUI 配置文件（可选）。 \\
\cmd{import eco <InputFileName>} &
导入先前的 eco 文件（可选）。 \\
\cmd{setup design} &
导入设计数据后设置数据库。 \\
\cmd{import apl <InputFileName>} &
导入 APL 生成的 \texttt{<cell>.current} 文件（动态）。 \\
\cmd{import apl -c <InputFileName>} &
导入 APL 生成的 \texttt{<cell>.cdev} 文件（动态）。 \\
\cmd{import avm <configFileName>} &
由配置文件生成并导入存储器模型、\texttt{vmemory.current} 与 \texttt{vmemory.cdev} 文件。 \\
\cmd{import db <InputFileName>} &
导入先前在 setup、提取或仿真后生成的结果。 \\
\bottomrule
\end{tabularx}
\end{table}

现在可以加载工艺与设计信息，并运行 RedHawk 静态 IR drop 分析。

\begin{enumerate}
  \item 参见附录~\ref{chap:install}「设置 RedHawk 环境」，了解设置许可、环境变量与路径以访问 RedHawk 二进制文件的流程。
  \item 在 UNIX shell 中输入以下命令：
\begin{lstlisting}
redhawk
\end{lstlisting}
  该命令在屏幕上显示执行消息；以下命令将执行消息保存到 \texttt{apache.log}：
\begin{lstlisting}
redhawk >& apache.log&
\end{lstlisting}
  RedHawk GUI 将显示。

  \item 若使用命令行准备设计数据，按上一节及附录~\ref{chap:files}「文件定义」编译所需输入文件。

  \item 若所有所需输入文件均在 GSR 文件中引用或包含，则使用：
\begin{lstlisting}
import gsr <filename>
\end{lstlisting}

  \item 若并非所有所需文件均在 GSR 中引用，则使用上表中的 TCL 命令导入所需文件，例如 LEF、DEF、LIB 与 \texttt{*.tech}。

  \item 文件导入后，要设置设计数据库，输入：
\begin{lstlisting}
setup design
\end{lstlisting}
  以下数据完整性检查确保 RedHawk 数据适当：
  \begin{enumerate}
    \item 频率不超过 50\,GHz。
    \item 转换时间不超过 5\,ns。
    \item 动态压降分析的仿真时间或帧大小不超过 200\,ns；上电分析不超过 2\,\textmu s。
    \item SPICE subcircuit 预解析确保无语法或 include 错误，且与 \texttt{.ploc} 无未匹配节点。
    \item switch model 预解析确保无语法错误且数值合理。
  \end{enumerate}
\end{enumerate}

\subsubsection{GUI 数据加载}
\paragraph*{GUI Data Loading}

\begin{enumerate}
  \item 参见附录~\ref{chap:install}「设置 RedHawk 环境」，了解设置许可、环境变量与路径以访问 RedHawk 二进制文件的流程。
  \item 若使用 GUI 准备设计数据，先从命令行启动 RedHawk 以显示界面：
\begin{lstlisting}
redhawk
\end{lstlisting}
  GUI 将显示。
  \item 选择菜单命令 \textbf{File $\rightarrow$ Import Design Data}。

  将显示「Import Design Data」对话框。
  \item 对每个所需输入文件，点击 \textbf{Select Filename} 按钮，在显示的对话框中输入路径与文件名。

  若所有输入文件与 \texttt{*.gsr} 文件路径相同，可点击表单底部的 \textbf{Use Same Prefix} 按钮，自动为所有输入文件填入相同路径。
  \item 完成所有输入文件的路径与文件名。
  \item 确保 \textbf{Database setup} 按钮处于按下状态。
  \item 点击 \textbf{OK} 加载所有文件并设置设计数据库。

  加载 DEF 文件时，可使用仅 VDD 或仅 VSS 的 DEF，以加速仅 VDD 或仅 VSS 分析。
  \item 所有文件加载后，查看出现的版图视图：
  \begin{itemize}
    \item 使用鼠标右键缩放。使用显示区右侧的 \textbf{View} 按钮更改其他视图方面。
    \item 使用右下角的 Mini-view 面板导航。
    \item 使用右侧 \textbf{Configuration} 面板中的 \textbf{View Layers} 按钮开关层显示。
    \item 使用鼠标左键选择设计对象。选择后，对象以黄色高亮显示。
  \end{itemize}
\end{enumerate}

\begin{seeAlsoBox}
继续分析流程见第~\ref{chap:static}~章「功耗计算、静态 IR Drop 与 EM 分析」。
\end{seeAlsoBox}

\subsection{RedHawk 配置文件}
\subsubsection*{RedHawk Configuration Files}

除输入数据文件外，部分 RedHawk 功能与工具还需要必须准备的配置文件，如表~\ref{tab:rh-config} 所列。

\begin{table}[htbp]
\centering
\caption{RedHawk 配置文件（表 3-5）}
\label{tab:rh-config}
\small
\begin{tabularx}{\textwidth}{@{}lX@{}}
\toprule
文件名 & 用途 \\
\midrule
\texttt{apl\_config\_file} &
Apache Power Library（APL）程序用于表征单元、decap、I/O 与存储器，以进行动态压降分析。见第~\ref{chap:apl}~章 APL 表征。 \\
\texttt{switch\_config\_file} &
\texttt{aplsw} 工具用于表征 header/footer 开关，和/或为低功耗设计中使用的开关生成分段线性（PWL）电流模型。见第~\ref{chap:lp}~章低功耗分析。 \\
\texttt{avm\_config\_file} &
\texttt{avm}（\cmd{import avm}）工具用于为存储器生成开关电流剖面。 \\
\texttt{sim2iprof\_config\_file} &
\texttt{sim2iprof} 工具用于从第三方仿真输出中提取 Read/Write/Standby 信号，并生成电流剖面波形。见附录~\ref{chap:util}。 \\
\texttt{gds2rh\_config\_file} &
\texttt{gds2rh} 工具用于将 GDSII 文件转换为 DEF 格式。见附录~\ref{chap:util}。

\texttt{gds2rh -m} 工具用于生成存储器及其他 IP 的详细视图。见附录~\ref{chap:util}。 \\
\texttt{psi.ctl} &
PJX 配置与控制文件，支持 PJX 时序与时钟树分析程序。见第~\ref{chap:noise}~章。 \\
\bottomrule
\end{tabularx}
\end{table}
